\section*{Full outline}
\addcontentsline{toc}{section}{Full outline}

This volume gives beginners the theory they can hear, write, and judge, and ties every concept back to Yeban's piano roll, arrangement area, and mixer. It is not a compressed conservatory textbook but a self-study route: identify pitch, count rhythm, hear intervals, and finally write an eight-bar passage with functional harmony.

\subsection*{Audience and prerequisites}

\begin{itemize}
  \item No score-reading or instrument skills required; you only need to open Yeban and use the piano-roll draw tool.
  \item Every chapter starts with ``what you can do afterwards'' and ends with ``how to do it in Yeban'' plus exercises; answers appear at the end of the chapter.
  \item A term is introduced in Chinese (English) at first use; the full bilingual list is Appendix A. The body avoids synonym piles.
\end{itemize}

\subsection*{Twelve-chapter route}

\begin{enumerate}[label=\textbf{Chapter \arabic*}]
  \item \textbf{Sound and pitch}\\
  Vibration and frequency, note names, octaves, scientific pitch notation, MIDI numbers, 12-tone equal temperament, middle C, and A4 = 440 Hz; align the piano keyboard, Yeban's pitch rows, and the frequency formula.
  \item \textbf{Rhythm and meter}\\
  Pulse and beat, time signatures (2/4, 3/4, 4/4, 6/8), note values and rests, dots and ties, syncopation and backbeat, tempo and BPM; count cells on Yeban's timeline and align to the transport.
  \item \textbf{Intervals}\\
  Semitones and whole tones, major/minor/augmented/diminished intervals, simple and compound intervals, consonance and dissonance, interval recognition and singing; measure the row distance between two notes in the piano roll.
  \item \textbf{Scales and modes}\\
  Major and minor, key signatures and accidentals, pentatonic scales, common modes (Dorian, Lydian, Mixolydian), relative major/minor; check scale degrees with Yeban's scale/chord helpers.
  \item \textbf{Chords}\\
  Triads, seventh chords, inversions, chord symbols (root + quality + inversion), common progressions (I–V–vi–IV, ii–V–I); stack notes in the roll and listen to voicing and low-end weight in the mixer.
  \item \textbf{Harmonic progression and function}\\
  Tonic/dominant/subdominant function, cadences (authentic, plagal, half, deceptive), borrowed chords, secondary dominants, introduction to modulation; explain ``why this chord follows that one'' as operable rules.
  \item \textbf{Melody writing}\\
  Motif, phrase, call and response, repetition and sequence, melody versus chord tones and tension notes, rhythmic motives; write a four-bar call-and-response in Yeban and vary it twice.
  \item \textbf{Introduction to form}\\
  Sections and layers, repetition and variation, common pop/folk forms (AABA, verse–chorus, intro–verse–chorus–bridge–outro), transitions; organize structure with clip duplication and naming in the arrangement area.
  \item \textbf{Score reading and notation}\\
  Staff basics, clefs, note and rest symbols, expression and dynamic markings, the role of numbered notation and TAB; map staff positions to piano-roll rows and read simple examples.
  \item \textbf{Blues, pentatonic and modal writing}\\
  Pentatonic as a pitch set, blues scale and blue notes, the 12-bar blues skeleton and substitutions, modal characteristic tones and modal vamps; turn ``two notes missing'' into a style rather than a simplified major scale.
  \item \textbf{Advanced harmony: sevenths, voice leading and reharmonisation}\\
  Five seventh-chord qualities, resolving the seventh and leading tone, parallel fifths/octaves, inversions and bass lines, drop-2, functional substitution, pedal bass and sequential reharmonisation; from ``can add chords'' to ``can change the harmony''.
  \item \textbf{Rhythm, phrase and form: advanced}\\
  6/8 versus 3/4, hemiola, three-against-two, syncopation and space, four motivic developments, phrase peak and cadence, section jobs and 16-bar structure; treat rhythm as a structural tool, not just a drum pattern.
\end{enumerate}

\subsection*{Fixed chapter structure}

\begin{description}[leftmargin=2em,style=nextline]
  \item[Concept] Start from everyday experience, give intuition first, then definitions and formulas.
  \item[Score/figure] Keyboard, roll grid, staff, or numbered-notation diagrams drawn in TikZ; no external images.
  \item[How to do it in Yeban] Only entries that actually exist in the live tip and the in-app manual; unavailable features are explicitly marked ``not yet available.''
  \item[Exercises] Four to six items with answers or hints; at least one must be done inside Yeban.
  \item[Help topic map] Section $\rightarrow$ in-app help anchor (English lowercase with hyphens).
\end{description}

\subsection*{Suggested pace}

\begin{itemize}
  \item Complete beginner: one chapter per week; read the concept, then do the exercises in Yeban. Do not skip Chapters 1–5.
  \item Instrumentalist new to digital audio: skim Chapters 1, 2, and 9; spend time on the writing exercises in Chapters 5–7.
  \item ``I just want to write a song in Yeban'': follow 1 $\rightarrow$ 2 $\rightarrow$ 5 $\rightarrow$ 7 $\rightarrow$ 8; use Chapters 3, 4, and 6 as reference.
\end{itemize}

\subsection*{Boundaries with the other volumes}

\begin{itemize}
  \item This volume explains why music works, not every Yeban menu; day-to-day UI steps belong to the Yeban User Guide.
  \item Instrument timbre, synth parameters, and effect chains belong to the Instruments and Synthesizers volume; this volume only touches orchestration and mixing where needed.
  \item Production workflow (track layering, buses, sends, mastering) belongs to the Arrangement and Orchestration volume; Chapters 6–12 here stay at the writing level.
\end{itemize}
